\newpage
\section{\tkzClass{regular polygon} 类}

变量 \tkzVar{regular\_polygon}{RP} 持有一个表，用于存放正多边形。它是可选的，你可以自由选择变量名。不过，为清晰和一致起见，推荐使用 \code{RP}。如果你使用自定义变量（例如 REP），则必须手动初始化它。如果使用了 \code{RP}，\tkzFct{tkz-elements}{init\_elements()} 函数会重新初始化该表。


\subsection{创建正多边形}

\tkzClass{regular\_polygon} 类根据圆心、第一个顶点和边数来构造一个正多边形。

\medskip
结果通常存放在 \tkzVar{regular\_polygon}{RP} 中。

\begin{mybox}
\begin{verbatim}
RP.hex = regular_polygon:new(z.O, z.A, 6)
\end{verbatim}
\end{mybox}

\textbf{简写形式。}

为简洁起见，使用简写形式：

\begin{mybox}
\begin{verbatim}
RP.hex = regular_polygon(z.O, z.A, 6)
\end{verbatim}
\end{mybox}

\subsection{Regular\_polygon 的属性}
\label{sub:regular_polygon_attributes}

\begin{mybox}
创建 | RP.IA = regular_polygon(z.I,z.A,6)|
\end{mybox}


  \bgroup
  \catcode`_=12
  \small
  \captionof{table}{Regular\_polygon 的属性。}\label{regular:attributes}
  \begin{tabular}{ll}
  \toprule
  \texttt{属性}      & \textbf{用法}  \\
 \midrule
  \tkzAttr{regular\_polygon}{center}  & |z.I = RP.IA.center|  \\
  \tkzAttr{regular\_polygon}{vertices}   & 包含所有顶点附标的数组  \\
  \tkzAttr{regular\_polygon}{through} & 第一个顶点  \\
  \tkzAttr{regular\_polygon}{circle}  & 定义以 I 为圆心、过 A 的圆 \\
  \tkzAttr{regular\_polygon}{type}    & |RP.IA.type= 'regular\_polygon'|   \\
  \tkzAttr{regular\_polygon}{side}    & |s = RP.IA.side| ; s = 边长\\
  \tkzAttr{regular\_polygon}{circumradius}&  |RP.IA.circumradius| ; 外接圆半径 \\
  \tkzAttr{regular\_polygon}{inradius}&  |RP.IA.inradius| ; 内切圆半径   \\
  \tkzAttr{regular\_polygon}{apothem}    &  |RP.IA.apothem| ; 边心距（与内切圆半径相同）   \\
  \tkzAttr{regular\_polygon}{apothem\_foot} & |z.H = RP.IA.apothem_foot| ; 圆心在第一条边上的投影 \\
  \tkzAttr{regular\_polygon}{nb}      & |RP.IA.nb| ; 边数 \\
  \tkzAttr{regular\_polygon}{perimeter} & |RP.IA.perimeter| ; 周长 \\
  \tkzAttr{regular\_polygon}{area}    & |RP.IA.area| ; 面积 \\
  \tkzAttr{regular\_polygon}{angle}   &  |RP.IA.angle| ; 圆心与两个相邻顶点所成的角   \\
  \bottomrule %
  \end{tabular}
  \egroup




\subsubsection{五边形}
\label{ssub:pentagon}
\begin{minipage}{.5\textwidth}
\begin{verbatim}
\directlua{
init_elements()
z.O = point(0, 0)
z.I = point(1, 3)
z.A = point(2, 0)
RP.five  = regular_polygon(z.I, z.A, 5)
RP.five:name("P_")
C.ins  = circle:radius(z.I,
             RP.five.inradius)
z.H = RP.five.apothem
}
\begin{tikzpicture}[scale = .75]
\def\nb{\tkzUseLua{RP.five.nb}}
\tkzGetNodes
\tkzDrawCircles(I,A I,H)
\tkzDrawPolygon(P_1,P_...,P_\nb)
\tkzDrawPoints[red](P_1,P_...,P_\nb,H,I)
\tkzLabelPoints[red](I,A,H)
\end{tikzpicture}
\end{verbatim}
\end{minipage}
\begin{minipage}{.5\textwidth}
\directlua{
   init_elements()
   z.O = point(0, 0)
   z.I = point(1, 3)
   z.A = point(2, 0)
   RP.five = regular_polygon(z.I, z.A, 5)
   RP.five:name("P_")
   C.ins  = circle:radius(z.I, RP.five.inradius )
   z.H = RP.five.apothem_foot
  }

\begin{center}
  \begin{tikzpicture}[scale = .75]
  \def\nb{\tkzUseLua{RP.five.nb}}
  \tkzGetNodes
  \tkzDrawCircles(I,A I,H)
  \tkzDrawPolygon(P_1,P_...,P_\nb)
  \tkzDrawPoints[red](P_1,P_...,P_\nb,H,I)
  \tkzLabelPoints[red](I,A,H)
  \end{tikzpicture}
\end{center}


\end{minipage}

\subsection{Regular\_polygon 的方法}
\label{sub:regular_polygon_methods}


\bgroup
\catcode`_=12
\small
\captionof{table}{regular\_polygon 的方法。}\label{regular:methods}
\begin{tabular}{ll}
\toprule
\texttt{构造器}&\\  \\
\midrule
\tkzFct{regular\_polygon}{new(O,A,n)} & [\ref{ssub:function_regular_new}] \\
\midrule
\texttt{返回圆的方法}&\\
\midrule
\tkzMeth{regular\_polygon}{incircle()} & [\ref{ssub:method_regular_incircle}]\\
\midrule
\texttt{返回点的方法}&\\
\midrule
\tkzMeth{regular\_polygon}{get(<i>)} & [\ref{ssub:method_regular_get}] \\
\midrule
\texttt{其他方法}&\\
\midrule
\tkzMeth{regular\_polygon}{name (string)} & [\ref{ssub:method_regular_name}] \\
\bottomrule %
\end{tabular}
\egroup


\subsubsection{方法 \tkzMeth{regular\_polygon}{new(pt, pt, n)}}
\label{ssub:function_regular_new}

\begin{tkzexample}[latex=.5\textwidth]
\directlua{
  z.A = point(0, -4)
  z.O = point(0, 0)
  RP.six = regular_polygon(z.O, z.A, 6)
  RP.six:name("A_")
  z.i, z.p = RP.six:incircle():get()}

\begin{tikzpicture}[gridded,scale=.75]
  \tkzGetNodes
  \tkzDrawCircles[red](O,A)
  \tkzDrawCircles[teal](i,p)
  \tkzDrawPolygon(A_1,A_...,A_6)
  \tkzDrawPoints[red](A_1,A_...,A_6)
  \tkzLabelPoints[red](A_1,A_...,A_6)
\end{tikzpicture}
\end{tkzexample}

\subsubsection{方法 \tkzMeth{regular\_polygon}{get(<i>)}}
\label{ssub:method_regular_get}

\code{get()} 返回所有顶点；\code{get(i)} 返回第 i 个顶点（$1 \le i \le n$）。无效的索引会引发错误。

\begin{verbatim}
z.A, z.B, z.C = RP.tri:get()
z.B = RP.tri:get(2)
\end{verbatim}

\subsubsection{方法 \tkzMeth{regular\_polygon}{incircle()}}
\label{ssub:method_regular_incircle}

见前一个示例 [\ref{ssub:function_regular_new}]

\subsubsection{方法 \tkzMeth{regular\_polygon}{name(s)}}
\label{ssub:method_regular_name}

\code{RP.six:name("A\_")} 会创建点 \code{z.A\_1}、……、\code{z.A\_6}。

见 [\ref{ssub:function_regular_new}]
\endinput